\section*{Hoofdstuk \ref{ch:g7:identifying-substances} --- Stoffen herkennen}

\begin{solution}{exo:g7:identifying-substances:1}
De kalkwatertest: je leidt het gas door helder kalkwater, dat troebel
wordt als het gas koolstofdioxide is.
\end{solution}

\begin{solution}{exo:g7:identifying-substances:2}
Het poeder is watervrij kopersulfaat, en de vloeistof bevat water.
\end{solution}

\begin{solution}{exo:g7:identifying-substances:3}
Zuurstof.
\end{solution}

\begin{solution}{exo:g7:identifying-substances:4}
Potlood lost niet op in het \omterm{def:g6:solutions-and-solubility:solute}{oplosmiddel}. Een inktlijn zou zelf
meegenomen en gescheiden worden, en het \omterm{def:g7:identifying-substances:chromatography}{chromatogram} bederven.
\end{solution}

\begin{solution}{exo:g7:identifying-substances:5}
Als de stip onder het \omterm{def:g6:solutions-and-solubility:solute}{oplosmiddel} lag, zouden de kleurstoffen in het
\omterm{def:g6:solutions-and-solubility:solute}{oplosmiddel} in het bekerglas \omterm{def:g2:mixing-and-dissolving:dissolve}{oplossen} in plaats van in het papier omhoog
te kruipen.
\end{solution}

\begin{solution}{exo:g7:identifying-substances:6}
Nee, het is een \omterm{def:g6:pure-substances-and-mixtures:mixture}{mengsel} van twee kleurstoffen: waarschijnlijk de gele en
de blauwe kleurstof van de referenties.
\end{solution}

\begin{solution}{exo:g7:identifying-substances:7}
Vang een beetje van het gas op in een reageerbuis en houd een brandende
spaander bij de opening: een blaffend plofje wijst op waterstof.
\end{solution}

\begin{solution}{exo:g7:identifying-substances:8}
Uitgeademde \omterm{def:g4:air-a-mixture-of-gases:air}{lucht} bevat veel meer koolstofdioxide dan frisse \omterm{def:g4:air-a-mixture-of-gases:air}{lucht}.
\end{solution}

\begin{solution}{exo:g7:identifying-substances:9}
De test laat zien dat er water aanwezig is, niet dat het er alleen is:
zout water, sap of elk ander waterig \omterm{def:g6:pure-substances-and-mixtures:mixture}{mengsel} zou het poeder ook blauw
kleuren.
\end{solution}

\begin{solution}{exo:g7:identifying-substances:10}
Drie kleurstoffen (drie vlekken). De gele kleurstof is het hoogst
gekomen.
\end{solution}

\begin{solution}{exo:g7:identifying-substances:11}
Bijvoorbeeld: een gloeiende spaander in elke pot --- de pot waarin hij
weer opvlamt, bevat zuurstof. Bij de andere twee een brandende spaander
bij de opening --- een blaffend plofje wijst op waterstof. De laatste pot
bevat koolstofdioxide (te bevestigen met kalkwater). Twee of drie tests
zijn genoeg.
\end{solution}

\begin{solution}{exo:g7:identifying-substances:12}
Het gas bevat misschien geen koolstofdioxide, of te weinig om in die tijd
iets te laten zien. Een controleproef met uitgeademde \omterm{def:g4:air-a-mixture-of-gases:air}{lucht}, waarvan je
weet dat ze koolstofdioxide bevat, controleert of het kalkwater werkt.
\end{solution}

\begin{solution}{pb:g7:identifying-substances:1}
\textbf{1.} Zodat de omstandigheden (papier, \omterm{def:g6:solutions-and-solubility:solute}{oplosmiddel}, tijd) voor alle
vier hetzelfde zijn: alleen dan kun je de hoogtes van de vlekken
vergelijken.

\textbf{2.} Een \omterm{def:g6:pure-substances-and-mixtures:mixture}{mengsel}: de inkt geeft drie aparte vlekken, drie
kleurstoffen.

\textbf{3.} Twee kleurstoffen. Nee: de inkt van het briefje heeft drie
kleurstoffen, en stift 1 heeft de rode en de gele kleurstof geen van
beide.

\textbf{4.} Haar drie vlekken liggen niet op dezelfde hoogtes als die
van het briefje: het zijn drie andere kleurstoffen.

\textbf{5.} Stift 2: haar drie vlekken liggen precies op de hoogtes van
de drie vlekken van het briefje.

\textbf{6.} Zuurstof.

\textbf{7.} Waterstof.

\textbf{8.} Koolstofdioxide.

\textbf{9.} Een blanco (een controleproef). Ze laat zien dat het
kalkwater niet vanzelf troebel wordt, en ook niet met gewone \omterm{def:g4:air-a-mixture-of-gases:air}{lucht} in
die tijd: de troebeling met gas Z komt dus echt van koolstofdioxide.

\textbf{10.} ``Water --- watervrij kopersulfaat --- wit wordt blauw'':
de vloeistof bevat water.

\textbf{11.} Nee. Een positieve test laat zien dat er water aanwezig is,
niet dat er niets anders is: de vloeistof kan een oplossing of een ander
waterig \omterm{def:g6:pure-substances-and-mixtures:mixture}{mengsel} zijn.

\textbf{12.} De inkt van het briefje bevat 3 kleurstoffen, en stift 2
heeft het briefje geschreven.
\end{solution}
